
\subsubsection{Architectures}

ResNet-18 with Batch Normalization (ResNet-18-BN) uses standard torchvision ResNet-18 with 17 Batch Normalization layers. Batch Normalization normalizes activations per channel using batch statistics:

$$
\hat{x}_i = \frac{x_i - \mu_B}{\sqrt{\sigma_B^2 + \epsilon}}, \quad y_i = \gamma \hat{x}_i + \beta
$$

where $\mu _B$ and $\sigma ^{2}_B$ are mean and variance computed over the batch dimension.

ResNet-18 with Layer Normalization (ResNet-18-LN) replaces all Batch Normalization layers with Layer Normalization. Layer Normalization normalizes per-instance over feature dimensions:

$$
\hat{x}_i = \frac{x_i - \mu_L}{\sqrt{\sigma_L^2 + \epsilon}}, \quad y_i = \gamma \hat{x}_i + \beta
$$

where $\mu _L$ and $\sigma ^{2}_L$ are mean and variance computed over [C, H, W] for each instance. Layer Normalization processes each example independently, eliminating batch-level information aggregation.

\subsubsection{Accuracy-Matched Comparison}

Batch Normalization accelerates training convergence \citep{Santurkar2019How}. Comparing at fixed epochs confounds architectural effects with training speed. We compare at matched average accuracy checkpoints. For each seed:

\begin{enumerate}
\item Train ResNet-18-BN and ResNet-18-LN for up to 100 epochs.
\item Identify the epoch where each architecture first reaches 90\% average accuracy.
\item Record worst-group accuracy at those epochs.
\item Compute worst-group gap = average accuracy - worst-group accuracy.
\item Perform paired t-test across 10 seeds.
\end{enumerate}

\subsubsection{Gradient Measurement}

To establish a mechanistic explanation, we measure gradient flow during early training (epochs 0--19). We instrument conv1.weight with gradient hooks:

\begin{enumerate}
\item Partition training samples into majority groups (spurious-aligned) and minority groups (spurious-misaligned).
\item Compute gradient contributions from each group: g\_maj and g\_min.
\item Compute gradient ratio r = ||g\_maj||$_2$ / ||g\_min||$_2$.
\item Average over epochs 0--19 for each seed.
\item Compare across architectures with paired t-test.
\end{enumerate}

\subsubsection{Training Configuration}

All experiments use identical hyperparameters:
\begin{itemize}
\item Optimizer: SGD (lr=0.01, momentum=0.9, weight\_decay=1e-4)
\item Learning rate: constant (no schedule)
\item Batch size: 64
\item Initialization: He normal
\item Epochs: 20 (proof-of-concept)
\item Seeds: 10 (0--9)
\end{itemize}

\subsubsection{Dataset}

Planned: Waterbirds dataset \citep{Sagawa2020Distributionally} with 4795 training images, 95\% spurious correlation (landbird-land background, waterbird-water background).

Used: Synthetic spurious correlation dataset with 5000 training samples, 1000 test samples, 90\% spurious correlation. Generated due to WILDS Waterbirds server unavailability (HTTP 500 error).

\subsubsection{Evaluation Metrics}

\begin{itemize}
\item Average accuracy: accuracy over all test samples
\item Worst-group accuracy (WGA): minimum accuracy across four groups (label $\times$ spurious feature)
\item Worst-group gap: average accuracy - WGA
\end{itemize}
